\documentclass[11pt,a4paper]{article}

\usepackage{fontspec}
\IfFontExistsTF{Times New Roman}{\setmainfont{Times New Roman}}{\setmainfont{Tinos}}
\IfFontExistsTF{Consolas}{\setmonofont{Consolas}[Scale=0.92]}{\setmonofont{Liberation Mono}[Scale=0.92]}
\usepackage[english]{babel}
\usepackage{geometry}
\usepackage{xcolor}
\usepackage{enumitem}
\usepackage{listings}
\usepackage{booktabs}
\usepackage{array}
\usepackage{amsmath,amssymb}
\usepackage{titlesec}
\usepackage{fancyhdr}
\usepackage{parskip}
\usepackage{tcolorbox}
\tcbuselibrary{skins}
\usepackage{tabularx}
\usepackage{lastpage}
\usepackage{hyperref}

\geometry{top=1.5cm,bottom=1.5cm,left=1.7cm,right=1.7cm}

\definecolor{ForestGreen}{RGB}{22,100,70}
\definecolor{SoftGreen}{RGB}{229,244,235}
\definecolor{CodeBg}{RGB}{248,250,249}
\definecolor{DarkGray}{RGB}{55,55,55}
\definecolor{primaryblue}{RGB}{22,100,70}
\definecolor{secondaryteal}{RGB}{22,100,70}
\definecolor{boxbg}{RGB}{248,250,249}
\definecolor{borderblue}{RGB}{159,198,178}

\setlength{\parindent}{0pt}
\setlength{\parskip}{0pt}
\setlist[itemize]{leftmargin=1.3em,itemsep=1pt,topsep=1pt,parsep=0pt,partopsep=0pt}
\setlist[enumerate]{leftmargin=1.5em,itemsep=1.5pt,topsep=1pt,parsep=0pt,partopsep=0pt}
\setlist[enumerate,2]{topsep=0.5pt,itemsep=0pt,parsep=0pt,partopsep=0pt}

\lstset{
  language=Python,
  basicstyle=\ttfamily\scriptsize,
  keywordstyle=\color{blue!70!black}\bfseries,
  stringstyle=\color{red!65!black},
  commentstyle=\color{ForestGreen},
  showstringspaces=false,
  columns=fullflexible,
  keepspaces=true,
  breaklines=true,
  frame=single,
  rulecolor=\color{ForestGreen!45},
  backgroundcolor=\color{CodeBg},
  xleftmargin=0.2em,
  xrightmargin=0.2em,
  aboveskip=0.2em,
  belowskip=0.2em,
  tabsize=4
}

\pagestyle{fancy}
\fancyhf{}
\setlength{\headheight}{14pt}
\fancyhead[L]{\footnotesize\textbf{DSAI1003 -- Fundamental Programming Concepts in Python}}
\fancyhead[R]{\footnotesize\textbf{Numbers, Strings \& Operators Exam -- 60 Minutes}}
\fancyfoot[L]{\scriptsize Faculty of Data Science and Artificial Intelligence -- National Economics University}
\fancyfoot[R]{\footnotesize Page \thepage\ of 6}
\renewcommand{\headrulewidth}{0.4pt}
\renewcommand{\footrulewidth}{0.3pt}

\titleformat{\section}{\large\bfseries\color{ForestGreen}}{}{0pt}{}
\titleformat{\subsection}{\normalsize\bfseries\color{ForestGreen}}{}{0pt}{}
\titlespacing*{\section}{0pt}{4pt}{2pt}
\titlespacing*{\subsection}{0pt}{3pt}{1.5pt}

\newtcolorbox{headerbox}{
  colback=white,
  colframe=ForestGreen,
  boxrule=0.7pt,
  arc=1.5mm,
  left=2.5mm,
  right=2.5mm,
  top=1.5mm,
  bottom=1.5mm
}

\newtcolorbox{answerbox}[2][]{
  enhanced,
  colback=white,
  colframe=DarkGray!45,
  boxrule=0.5pt,
  arc=1mm,
  left=2mm,
  right=2mm,
  top=1.2mm,
  bottom=1.2mm,
  title=\footnotesize\color{DarkGray}\textbf{#2},
  #1
}

\hypersetup{
  colorlinks=true,
  linkcolor=ForestGreen,
  urlcolor=ForestGreen,
  pdftitle={DSAI1003 - Numbers, Strings, and Operators in Python Exam},
  pdfauthor={Faculty of Data Science and Artificial Intelligence - National Economics University (NEU)}
}

\begin{document}

% =========================================================================
% PAGE 1: OFFICIAL HEADER + ANSWER SHEET + SECTION A (Q1 - Q5)
% =========================================================================

\noindent
\begin{minipage}[t]{0.64\textwidth}
  \raggedright\small
  \textbf{NATIONAL ECONOMICS UNIVERSITY}\\
  \textbf{COLLEGE OF TECHNOLOGY}\\
  \textbf{FACULTY OF DATA SCIENCE \& ARTIFICIAL INTELLIGENCE (FDA)}
\end{minipage}%
\hfill
\begin{minipage}[t]{0.34\textwidth}
  \raggedleft\footnotesize
  \textbf{OFFICIAL MIDTERM EXAMINATION}\\
  Academic Year: 2026--2027\\
  Semester: Fall 2026 \quad Course Code: \textbf{DSAI1003}
\end{minipage}

\vspace{0.2em}

\begin{center}
  {\Large\bfseries\color{ForestGreen} FUNDAMENTAL PROGRAMMING CONCEPTS IN PYTHON}\\[0.10em]
  {\normalsize\bfseries PROGRESS EXAM: NUMBERS, STRINGS, AND OPERATORS}\\[0.10em]
  {\normalsize\textbf{Paper Code:} \rule{1.6cm}{0.4pt} \quad | \quad \textbf{Duration: 60 Minutes} (Excluding paper distribution)}
\end{center}

\vspace{0.1em}

\begin{headerbox}
\begin{tabular*}{\textwidth}{@{\extracolsep{\fill}}ll@{}}
  \textbf{Candidate Full Name:} \dotfill & \textbf{Student ID (MSSV):} \dotfill \\[0.25em]
  \textbf{Class / Cohort:} \dotfill & \textbf{Exam Room / Seat No.:} \dotfill
\end{tabular*}

\vspace{0.25em}
\renewcommand{\arraystretch}{1.1}
\setlength{\tabcolsep}{4.5pt}
\centering
\begin{tabular}{|l|c|c|c|c||c|c|}
\hline
\textbf{Section} & \textbf{A} & \textbf{B} & \textbf{C} & \textbf{D} & \textbf{Total (100)} & \textbf{Examiner Signature} \\
\hline
\textbf{Max Score} & 20 & 25 & 25 & 30 & \textbf{100} & \\[0.35em]
\cline{1-6}
\textbf{Score} & & & & & & \\[0.45em]
\hline
\end{tabular}

\vspace{0.2em}
\raggedright
\scriptsize
\textbf{Instructions:} Closed-book examination. All electronic devices and internet access are prohibited. Answer all questions directly in the designated spaces. In programming tasks, use the \textbf{Input $\rightarrow$ Process $\rightarrow$ Output} pattern. Do \textbf{not} define user functions (\texttt{def}) or use loops/branching. Assume users enter data in the required format. Do not detach pages.
\end{headerbox}

\vspace{1mm}
\noindent
\textbf{MULTIPLE-CHOICE ANSWER SHEET} \textit{(Write A, B, C, or D in the corresponding box)}:
\begin{center}
\setlength{\tabcolsep}{5pt}
\renewcommand{\arraystretch}{1.05}
\begin{tabular}{|c|c|c|c|c|c|c|c|c|c|c|}
\hline
\textbf{Question} & \textbf{1} & \textbf{2} & \textbf{3} & \textbf{4} & \textbf{5} & \textbf{6} & \textbf{7} & \textbf{8} & \textbf{9} & \textbf{10} \\
\hline
\textbf{Answer} & & & & & & & & & & \\[2mm]
\hline
\end{tabular}
\end{center}
\vspace{-2mm}
\section*{Section A: Numbers, Strings \& Operators Concepts \hfill \normalsize\normalfont(20 points)}
\textit{Answer all 10 questions. Each correct answer carries 2 points. Suggested time: 12 minutes.}

\begin{enumerate}[label=\textbf{Q\arabic*.},leftmargin=2em,itemsep=2pt,topsep=1pt]

\item Which of the following is a \textbf{valid} variable name according to Python identifier rules?
\begin{enumerate}[label=\Alph*.,leftmargin=2em]
  \item \texttt{2nd\_total}
  \item \texttt{tax-rate}
  \item \texttt{\_unit\_price}
  \item \texttt{class}
\end{enumerate}

\item What is the precise effect of executing the assignment statement \texttt{count = count + 1}?
\begin{enumerate}[label=\Alph*.,leftmargin=2em]
  \item It checks whether \texttt{count} is mathematically equal to \texttt{count + 1}.
  \item It evaluates \texttt{count + 1} using the current value of \texttt{count} and assigns the result back to \texttt{count}.
  \item It raises a syntax error because a variable cannot appear on both sides of \texttt{=}.
  \item It permanently converts the data type of \texttt{count} into a floating-point number.
\end{enumerate}

\item What value and data type are produced by the floor division expression \texttt{19 // 4} in Python?
\begin{enumerate}[label=\Alph*.,leftmargin=2em]
  \item \texttt{4.75} (type \texttt{float})
  \item \texttt{4} (type \texttt{int})
  \item \texttt{3} (type \texttt{int})
  \item \texttt{4.0} (type \texttt{float})
\end{enumerate}

\item Following standard operator precedence, what is the result of the expression \texttt{2 + 3 * 2 ** 3}?
\begin{enumerate}[label=\Alph*.,leftmargin=2em]
  \item \texttt{26}
  \item \texttt{40}
  \item \texttt{200}
  \item \texttt{216}
\end{enumerate}

\item For any positive integer \texttt{n}, which expression correctly extracts its \textbf{last digit} (units digit)?
\begin{enumerate}[label=\Alph*.,leftmargin=2em]
  \item \texttt{n // 10}
  \item \texttt{n / 10}
  \item \texttt{n \% 10}
  \item \texttt{n \% 100}
\end{enumerate}

\end{enumerate}

\newpage

% ==============================================================================
% PAGE 2: SECTION A (Q6 - Q10) + SECTION B Q11
% ==============================================================================

\begin{enumerate}[label=\textbf{Q\arabic*.},leftmargin=2em,itemsep=2pt,topsep=1pt,resume]

\item Which statement regarding mathematical functions and the \texttt{math} module in Python is correct?
\begin{enumerate}[label=\Alph*.,leftmargin=2em]
  \item The \texttt{sqrt()} function is a built-in function that does not require an import statement.
  \item Built-in functions like \texttt{abs()} and \texttt{round()} can be called directly without importing \texttt{math}.
  \item The constant \texttt{math.pi} is a built-in keyword in Python available in all scopes.
  \item The \texttt{math} module cannot be used with floating-point numbers.
\end{enumerate}

\item Why can the comparison \texttt{0.1 + 0.2 == 0.3} evaluate to \texttt{False} in Python?
\begin{enumerate}[label=\Alph*.,leftmargin=2em]
  \item Python forbids addition between decimal floating-point values.
  \item Binary floating-point representation causes slight roundoff errors for certain decimal fractions.
  \item The \texttt{==} operator cannot compare floating-point operands.
  \item Python automatically converts \texttt{0.1 + 0.2} into an integer before comparison.
\end{enumerate}

\item Given the string assignment \texttt{s = "PYTHON"}, what are the values of \texttt{s[1]} and \texttt{s[-1]}?
\begin{enumerate}[label=\Alph*.,leftmargin=2em]
  \item \texttt{'P'} and \texttt{'N'}
  \item \texttt{'Y'} and \texttt{'N'}
  \item \texttt{'Y'} and \texttt{'O'}
  \item \texttt{'P'} and \texttt{'O'}
\end{enumerate}

\item What is the exact result of the string expression \texttt{"Data" + "5" * 2}?
\begin{enumerate}[label=\Alph*.,leftmargin=2em]
  \item \texttt{"Data10"}
  \item \texttt{"Data55"}
  \item \texttt{"Data5Data5"}
  \item A \texttt{TypeError} occurs because strings cannot be multiplied by integers.
\end{enumerate}

\item Consider the code: \texttt{word = "python"}; \texttt{word.upper()}; \texttt{print(word)}. What is printed?
\begin{enumerate}[label=\Alph*.,leftmargin=2em]
  \item \texttt{PYTHON}
  \item \texttt{python}
  \item \texttt{Python}
  \item An \texttt{AttributeError} occurs because strings do not possess an \texttt{upper()} method.
\end{enumerate}

\end{enumerate}

\section*{Section B: Output Tracing, Type Conversion \& Error Diagnosis \hfill \normalsize\normalfont(25 points)}
\textit{Answer all 5 questions. Each question carries 5 points. Suggested time: 15 minutes.}

\subsection*{Q11. Exact Output Tracing with Arithmetic \& Modulo (5 points)}
Predict the \textbf{exact screen output}, including spaces, punctuation, and line breaks:

\begin{lstlisting}
a = 23
b = 5
quotient = a // b
remainder = a % b
power_val = b ** 2 - a
print("q:", quotient, "r:", remainder, sep=" | ", end=";\n")
print(f"Check: {quotient * b + remainder == a}")
print(f"Power diff: {power_val}")
\end{lstlisting}

\begin{answerbox}[height=3.0cm]{Output Box for Q11}
\end{answerbox}

\newpage

% ==============================================================================
% PAGE 3: SECTION B, Q12--Q15
% ==============================================================================

\subsection*{Q12. String Operations, Indexing and Length Properties (5 points)}
Consider the following Python program:

\begin{lstlisting}
title = "Computer"
n = len(title)
first = title[0]
mid = title[n // 2]
last = title[-1]
combo = (first + mid + last).lower()
print(f"n = {n}")
print(f"Parts: {first}-{mid}-{last}")
print(f"Result: {combo * 2}")
\end{lstlisting}

\begin{itemize}[leftmargin=2em]
  \item Output of \texttt{print(f"n = \{n\}")}: \underline{\hspace{5cm}}
  \item Output of \texttt{print(f"Parts: \{first\}-\{mid\}-\{last\}")}: \underline{\hspace{5cm}}
  \item Output of \texttt{print(f"Result: \{combo * 2\}")}: \underline{\hspace{5cm}}
  \item Explain why integer floor division (\texttt{n // 2}) is required rather than \texttt{n / 2} for the index: \hrulefill
\end{itemize}

\subsection*{Q13. Diagnose and Correct Common Type and Operator Errors (5 points)}
For each case, state the error type / problem and write a corrected Python statement.

\renewcommand{\arraystretch}{1.35}
\begin{tabularx}{\textwidth}{|p{5.1cm}|X|X|}
\hline
\textbf{Code / Situation} & \textbf{Problem / Error} & \textbf{Correction / Valid Code} \\
\hline
\texttt{radius = input("Radius: ")}\\
\texttt{area = 3.14159 * radius ** 2} & & \\[4mm]
\hline
\texttt{s = "Python"}\\
\texttt{s[0] = "J"}\\
Goal: change first letter to \texttt{'J'} & & \\[4mm]
\hline
\texttt{total = 150}\\
\texttt{msg = "Total is: " + total} & & \\[4mm]
\hline
\end{tabularx}

\subsection*{Q14. Formula Evaluation with the \texttt{math} Module (5 points)}
Predict the exact output of the following geometry calculation:

\begin{lstlisting}
import math

a = 3.0
b = 4.0
hypotenuse = math.sqrt(a ** 2 + b ** 2)
area = 0.5 * a * b
print(f"Sides: {a:.1f}, {b:.1f}")
print(f"Hypotenuse: {hypotenuse:.2f}")
print(f"Triangle Area: {area:.1f}")
\end{lstlisting}

\begin{answerbox}[height=2.35cm]{Output Box for Q14}
\end{answerbox}

\subsection*{Q15. Replacing Magic Numbers with Named Constants (5 points)}
Consider the following soda volume calculation:

\begin{lstlisting}
cans = 6
bottles = 2
total_liters = cans * 0.355 + bottles * 2.0
\end{lstlisting}

\begin{itemize}[leftmargin=2em]
  \item State the two \textbf{magic numbers} present in line 3: \hrulefill
  \item What is the Python naming convention (PEP 8) for named constants? \hrulefill
  \item Rewrite line 3 using well-named constants \texttt{CAN\_VOLUME} and \texttt{BOTTLE\_VOLUME}:\\
  \rule{\linewidth}{0.4pt}
\end{itemize}

\newpage

% ==============================================================================
% PAGE 4: SECTION C
% ==============================================================================

\section*{Section C: Program Design Using the Input $\rightarrow$ Process $\rightarrow$ Output Model \hfill \normalsize\normalfont(25 points)}
\textit{Answer all tasks. Suggested time: 15 minutes.}

\begin{tcolorbox}[colback=boxbg,colframe=secondaryteal,arc=1.5mm,boxrule=0.6pt]
\textbf{Problem Specification: Cash Change Breakdown Calculator (Dollars, Quarters, Cents)}\\[1mm]
A cash register program receives an integer representing the total change due to a customer in \textbf{cents} (e.g., \texttt{387}). It computes the maximum number of full dollars (\$1 = 100 cents), the number of 25-cent quarters from the remainder, and the leftover cents. It also displays the total amount formatted as dollars and cents.
\end{tcolorbox}

\vspace{1mm}
\textbf{Task 1: Identify the IPO Components} (6 points)
\begin{itemize}[leftmargin=2em]
  \item \textbf{Input(s):} \hrulefill\\[4mm]
  \item \textbf{Process:} \hrulefill\\[4mm]
  \item \textbf{Output:} \hrulefill
\end{itemize}

\vspace{1mm}
\textbf{Task 2: Write the Program Logic in Ordered Steps} (7 points)\\
Write clear sequential steps using the Input $\rightarrow$ Process $\rightarrow$ Output structure. Your steps should state the required type conversion, integer floor division (\texttt{//}), modulo operations (\texttt{\%}), and output formatting.

\begin{answerbox}[height=8.0cm]{Program-Logic Work Area}
\end{answerbox}

\vspace{1mm}
\textbf{Task 3: Hand Trace} (6 points)\\
For \texttt{total\_cents = 387}, show the step-by-step arithmetic calculations and the exact formatted result.

\begin{answerbox}[height=3.7cm]{Hand-Trace Calculation}
\end{answerbox}

\vspace{1mm}
\textbf{Task 4: Explain Two Design Choices} (6 points)
\begin{enumerate}[label=\alph*),leftmargin=2em]
  \item Why must the user input be converted with \texttt{int()} rather than \texttt{float()} for this algorithm?\\[5mm]
  \hrulefill
  \item Why are floor division (\texttt{//}) and remainder (\texttt{\%}) used instead of true division (\texttt{/}) when finding coin quantities?\\[5mm]
  \hrulefill
\end{enumerate}

\newpage

% ==============================================================================
% PAGE 5: SECTION D, PROBLEM 1
% ==============================================================================

\section*{Section D: Applied Programming in Python \hfill \normalsize\normalfont(30 points)}
\textit{Write complete, syntactically correct Python statements using the \textbf{Input $\rightarrow$ Process $\rightarrow$ Output} pattern. Do \textbf{not} define custom functions (\texttt{def}), use loops, or use branching. Suggested time: 18 minutes.}

\subsection*{Problem 1: Cylindrical Storage Tank Volume and Coating Cost (15 points)}
Write a Python program that computes the geometry and protective coating cost for an industrial cylindrical storage tank.

\textbf{Requirements:}
\begin{enumerate}[leftmargin=2em,label=\arabic*.,itemsep=2pt]
  \item Prompt the user for the \texttt{tank\_id} (string) using \texttt{input()}.
  \item Prompt for the tank radius \texttt{radius} ($r$) in meters and convert it to \texttt{float}.
  \item Prompt for the tank height \texttt{height} ($h$) in meters and convert it to \texttt{float}.
  \item Import \texttt{math} to access \texttt{math.pi}. Define a constant \texttt{COATING\_COST\_PER\_SQM = 14.50}.
  \item Compute the storage volume: $V = \pi r^2 h$.
  \item Compute the total surface area (top, base, and side wall): $A = 2 \pi r h + 2 \pi r^2$.
  \item Compute the total coating cost: $\text{total\_cost} = A \times \text{COATING\_COST\_PER\_SQM}$.
  \item Display a formatted summary using f-strings: show \texttt{radius} and \texttt{height} with 2 decimal places, \texttt{volume} with 3 decimal places, \texttt{surface\_area} with 2 decimal places, and \texttt{total\_cost} with 2 decimal places prefixed by \texttt{\$}.
\end{enumerate}

The output should follow this general form:
\begin{lstlisting}
Tank ID: TANK-A1
Dimensions: Radius = 2.50 m, Height = 6.00 m
Storage Volume: 117.810 m^3
Total Surface Area: 133.52 m^2
Total Coating Cost: $1936.01
\end{lstlisting}

\begin{answerbox}[height=14.0cm]{Python Solution for Problem 1}
\end{answerbox}

\newpage

% ==============================================================================
% PAGE 6: SECTION D, PROBLEM 2
% ==============================================================================

\subsection*{Problem 2: Customer Voucher and Masked Account Statement (15 points)}
Write a Python program that processes customer data and formats a secure account voucher using string operations and numeric formatting.

\textbf{Requirements:}
\begin{enumerate}[leftmargin=2em,label=\arabic*.,itemsep=1.5pt]
  \item Prompt the user for the customer's \texttt{full\_name} (string) using \texttt{input()}.
  \item Prompt for an 8-character string \texttt{account\_code} (e.g., \texttt{ACCT8492}).
  \item Prompt for the current \texttt{account\_balance} in dollars and convert it to \texttt{float}.
  \item Define a named constant \texttt{DIVIDER\_LENGTH = 42}.
  \item Process the data:
  \begin{itemize}[itemsep=0.5pt,topsep=1pt]
    \item Convert \texttt{full\_name} to uppercase using \texttt{.upper()} and obtain its length using \texttt{len()}.
    \item Extract the first initial: first character of \texttt{full\_name} in uppercase.
    \item Construct a masked account code showing the first two characters, four asterisks (\texttt{"****"}), and the last two characters using string indexing, repetition, and concatenation.
    \item Create divider lines using string repetition: \texttt{"=" * DIVIDER\_LENGTH} and \texttt{"-" * DIVIDER\_LENGTH}.
  \end{itemize}
  \item Display the voucher formatted exactly as shown below, with \texttt{account\_balance} formatted to 2 decimal places.
\end{enumerate}

For example, if the user enters \texttt{Margaret Hamilton}, \texttt{ACCT8492}, and \texttt{1450.8}, the program outputs:
\begin{lstlisting}
==========================================
OFFICIAL CUSTOMER ACCOUNT VOUCHER
------------------------------------------
Customer Name: MARGARET HAMILTON (17 chars)
Initial & Masked Code: M | AC****92
Current Balance: $1450.80
==========================================
\end{lstlisting}

\begin{answerbox}[height=12.2cm]{Python Solution for Problem 2}
\end{answerbox}

\vfill
\begin{center}
  \rule{0.5\textwidth}{0.4pt}\\[1mm]
  \textbf{\color{primaryblue}--- END OF EXAM PAPER ---}
\end{center}

\end{document}
